\section{One finite cube for all parameters}\label{sec:unary-cube}

How many inputs does a counterexample require? A Dedekind packet uses
five naturals, of which four specify a polynomial. One natural already
suffices to exhibit a failure, even if its only interpretation is a
common search radius. The conversion in both directions has explicit
primitive recursive witnesses. The same construction applies to the
Dirichlet and ordinary packets.

Write \(e(\mathbf x)\) for the truncated excess in the relevant
packet. For RH the free coordinate is \(t\), for GRH the coordinates
are \((d_0,t)\), and for ERH they are \((d_0,H,u,v,t)\).
Define \(B_e(N)\) by scanning every tuple in
\(\{0,\ldots,N-1\}^r\), starting with zero and setting the flag to
one on a nonzero sample. If \(N=0\), the flag is zero for all three
families. For instance, \(B_e(2)\) checks the 32 combinations of
coordinates zero and one for ERH, including both polynomial words
in each combination. For GRH it checks four combinations. The cost
can be enormous because the underlying excess itself performs large
integer computations. The interest here lies in the arity and
logical strength of the test.

\begin{theorem}\label{thm:cubef}
Let \(A=\forall x_1\cdots\forall x_r\,[e(\mathbf x)=0]\), allowing
an arbitrary iPRA term and additional free parameters in \(e\).
The following are equivalent by derivations of object iPRA:
\[
 A,\qquad \forall N\ B_e(N)=0,\qquad
 \forall N\ B_e\bigl(N+g(N)\bigr)=0
\]
for each given unary primitive recursive program \(g\). In addition,
object iPRA proves the witness-preserving equivalence
\[
 \bigl(\exists\mathbf x\ e(\mathbf x)\ne0\bigr)
 \quad\longleftrightarrow\quad
 \bigl(\exists N\ B_e(N)\ne0\bigr).
\]
For a block of length zero the construction uses the term itself in
place of a flag.
\end{theorem}

\begin{proof}
The proof concerns formal derivations, so the induction formulas matter.
For a row of samples \(f(j,\mathbf p)\), define its scan by
\[
 S(0,\mathbf p)=0,\qquad
 S(j+1,\mathbf p)=
 \begin{cases}
    S(j,\mathbf p),&f(j,\mathbf p)=0,\\
    1,&f(j,\mathbf p)\ne0 .
 \end{cases}
\]
The conditional is a PR symbol with its two branches proved from the
recursion equations. The object theory proves
\(S(j+1)=0\leftrightarrow(f(j)=0\land S(j)=0)\):
induct on the sampled value in the quantifier-free formula for the
conditional. If all samples vanish, induct on the atomic formula
\(S(j)=0\) to obtain every clean prefix. Conversely, fix arbitrary
\(x\) and induct on the slack \(d\) in the quantifier-free implication
\[
 S\bigl(d+(x+1)\bigr)=0\ \longrightarrow\ f(x)=0 .
\]
At zero use the current-sample conjunct of the successor law.
At the successor use the previous-prefix conjunct and the induction
hypothesis. These are first-order derivations with all remaining
parameters free.

Compile a scoped term to a PR symbol by mapping each variable to a
projection and each application to composition. Structural recursion on
the finite term gives a formal proof of the compiler's defining
equality. At a leaf the cube code evaluates this term. At a universal
binder, apply the row scan to the child's cube code, keeping the same
horizon \(N\) inside each child and scanning the new coordinate only
up to \(N\). A permutation and a duplication of arguments implement
this convention with PR composition. This recursion runs over fixed
syntax while constructing the derivation. The only arithmetic
inductions within iPRA are quantifier-free.

For the forward direction, open the universal binder, prove the child
cube clean and use the row-scan result. For the converse take an
arbitrary tuple \(\mathbf x\). Choose
\[
                   N_{\mathbf x}=\sum_{i=1}^r(x_i+1).
\]
For each coordinate the sum of the other summands is an explicit
slack \(d_i\) satisfying \(N_{\mathbf x}=d_i+(x_i+1)\).
The same equations hold at horizon \(N_{\mathbf x}+g(N_{\mathbf x})\)
after increasing each slack by \(g(N_{\mathbf x})\).
The needed laws of addition have quantifier-free induction proofs.
Apply the single-row reading implication once per binder and conclude
\(e(\mathbf x)=0\). This covers every finite tuple with one and the
same horizon. Choosing the horizon after introducing the tuple
preserves the order of the object quantifiers. In particular, a known
bad tuple gives a bad cube of radius \(N_{\mathbf x}\).

For the converse on positive failures, one additional primitive
recursion keeps an index \(I\), initially zero. At stage \(j+1\),
retain \(I(j)\) if \(f(j)=0\), and set it to \(j\) otherwise.
The correctness invariant is the quantifier-free formula
\[
                 S(j)\ne0\ \longrightarrow\ f(I(j))\ne0 .
\]
At a zero sample any dirty flag comes from the preceding step.
At a nonzero sample the stored index points to the present failure.
Zero tests are decidable in iPRA, so the two branches suffice.
Starting with a bad cube, read this index in its outer scan. The child
cube at that index is bad. Repeat through the finite syntax down
to the nonzero excess. In each nonempty scan the hit is its last
failing index. This constructs the existential tuple directly.
\end{proof}

For example \(g(N)=N^2\) tests only sizes \(N^2+N\) and
\(g(N)=2^N\) tests only \(N+2^N\). Neither schedule has to enumerate
all sizes. Its explicit slack makes it cofinal in a way that iPRA can
use. The uniform Lean results are \texttt{CubeTests.diagonal\_iff},
\texttt{CubeTests.schedule\_iff}, and \texttt{CubeTests.failure\_iff}.
The specializations \texttt{rh\_cube}, \texttt{grh\_cube} and
\texttt{erh\_cube} have the original closed packet statements on their
left sides.

\paragraph{The first radius of failure.}
One may also count the number \(C(m)\) of dirty cubes with radius
strictly below \(m\). A counterexample to an entire family is
equivalent, with witnesses in both directions, to
\[
              \exists m\,[C(m)=0\ \land\ B_e(m)\ne0].
\]
Indeed a bad tuple supplies a dirty radius, and bounded search up to
that radius finds its first dirty cube. Its hit selectors then recover
a bad tuple in that cube. Here the minimum is taken across all field
codes and cutoffs at once. At the first dirty radius \(m\), every
bad tuple in that cube has some coordinate equal to \(m-1\), since its
predecessor cube is clean. In particular this witness minimizes
\(1+\max_i x_i\) across counterexamples. For the three nonempty blocks
the zero-radius cube is empty.

The Lean proof \texttt{failure\_iff\_first\_radius} composes
the positive equivalence above with the supplied one-dimensional
first-failure extractor. The companion \texttt{census\_iff}
transports that archive's full catalogue of tests to \(C\). For
example, each of \(\forall m\, C(m+1)=C(m)\) and
\(\forall m\,\exists d\, C(m+1+d)=0\) is an object-theoretically
equivalent formulation for each family.

